\documentclass{ximera}


\title{Math 1075 Midterm 3 Autumn 2019}

\begin{document}

\subsection*{Midterm Exam 3}
\subsection*{Autumn 2019}
\subsection*{Instructions:}
\begin{itemize}
  \item You have 55 minutes to complete this exam.
  \item Circle or box all final answers.
  \item A calculator may be used given the calculator policy outlined in the syllabus. Even if a calculator is used, you must show all work on each problem to receive full credit.
  \item Write clearly and legibly to receive credit.
  \item Do not round; give only exact answers (unless specified otherwise).
\end{itemize}

\begin{enumerate}
\item Graphs of two functions, $f(x)$ and $g(x)$, are given below. The graph of $f(x)$ is solid, and the graph of $g(x)$ is dashed. Use the graphs of the two functions to fill in the blanks below. (12 pts)\\
\includegraphics[width=\textwidth]{piecewise-functions-f-and-g}
\xmalt{Graph of two piecewise linear functions on the domain from $-1$ to $4$. The solid line $y = f(x)$ starts at $(-1, -1)$, rises to about $(0, 0.5)$, drops to $(0.5, 0)$, rises to $(1, 1)$, drops to $(1.5, 0.5)$, stays near $0.5$ through $x = 2$, then decreases to $(3, 0)$ and continues to $(4, 0)$. The dashed line $y = g(x)$ starts at $(-1, 1.75)$, decreases through $(0, 1.5)$ and $(1, 0.5)$, reaches a minimum near $x = 1.5$, then increases through $(3, 1)$ to $(4, 1.75)$.}

$$\begin{aligned}
& g(0)= \\
& (f+g)(-1)= \\
& (f-g)(1.5)=
\end{aligned}$$

$$\begin{aligned}
& (f g)(3)= \\
& (g \div f)(0)= \\
& (g \circ f)(0)=
\end{aligned}$$

\item Graphs of two relations, Relation A and Relation B, are given below. (8 pts)\\
\includegraphics[width=\textwidth]{relation-A-graph}
\xmalt{Relation A: A scatter plot with points at $(0, -1)$, $(1, -1)$, $(2, -1)$, $(3, -4)$, and $(4, 0)$ on a coordinate plane with $x$-axis from $0$ to $4$ and $y$-axis from $-4$ to $0$.}

\begin{figure}[h]
\begin{center}
  \includegraphics[width=\textwidth]{relation-B-graph}
\xmalt{Relation B: A piecewise linear function on a coordinate plane with $x$-axis from $0$ to $4$ and $y$-axis from $0$ to $-3$. A horizontal segment at $y = -1$ from $x = 1$ (open circle) to $x = 2$ (closed circle), then a line segment rising from $(2, -1)$ to $(3, -3)$ with an open circle at $(3, -3)$, and an isolated closed point at $(4, -1)$.}
\end{center}
\end{figure}

(a) Select the statement that is true.
\begin{itemize}
\item Relation A is a function, but Relation B is not.
\item Relation B is a function, but Relation A is not.
\item Both Relation A and Relation B are functions.
\item Neither Relation A nor Relation B are functions.
\end{itemize}
(b) Which of the following represents the domain of Relation A?
\begin{itemize}
\item $[1,2]$
\item $\{-1,1,2\}$
\item $[1,4]$
\item $\{0,1,2,3,4\}$
\end{itemize}
(c) Which of the following represents the range of Relation B?
\begin{itemize}
\item $\{1,2,3\}$
\item $[1,3)$
\item $\{1,2,3,4\}$
\item $(1,4)$
\end{itemize}
(d) Fill in the blank so that the ordered pair $(x, y)$ is included in Relation B: ( $\_\_\_\_$ 2). Select your answer from the options below.
\begin{itemize}
\item 1
\item 2
\item 3
\item 4
\end{itemize}

\item Simplify. Assume that the variable represents a positive real number. (8 pts)

$$\left(\frac{1}{32 x^{5}}\right)^{-2 / 5}$$

\item Simplify. Assume all variables represent positive real numbers. (10 pts)

$$\sqrt[4]{81 x^{19} z^{17}}$$

\item Simplify. Assume all variables represent positive real numbers. (8 pts)

$$x \sqrt{75 x v^{2}}+7 v \sqrt{27 x^{3}}$$

\item Solve for $y$, where $y$ is a real number. (11 pts)

$$y-4=\sqrt{-4 y+37}$$

\item Solve for $w$, where $w$ is a real number. (8 pts)

$$(w-7)^{2}-20=0$$

\item Simplify. Assume the variable represents a positive real number. (10 pts)

$$2 \sqrt[3]{500 x^{4}}-\sqrt[3]{256 x^{4}}$$

\item Simplify. Assume all variables represent positive real numbers. Be sure to rationalize the denominator. (9 pts)

$$\sqrt{\frac{7 h^{2} k^{4}}{32}}$$

\item Simplify. Assume the variable represents a positive real number. (8 pts)\\
$10 \sqrt[5]{x^{4}} \cdot 6 \sqrt{x}$

\item Rationalize the denominator and simplify. Assume that the variable represents a positive real number. (8 pts)\\
$\frac{-3}{3 \sqrt{w}+4}$

\end{enumerate}

\end{document}
